% Standalone review copy. To integrate into the thesis, retain the chapter body
% and merge the cited entries into the thesis bibliography. Required packages:
% amsmath, booktabs, longtable, array, natbib, hyperref.
% Numeric blocks are regenerated by build_chapter4_results.py.
% Figure boxes are author instructions, not claims that figures already exist.
\documentclass[12pt,a4paper]{report}
\usepackage[T1]{fontenc}
\usepackage[utf8]{inputenc}
\usepackage[margin=25mm]{geometry}
\usepackage{amsmath,booktabs,longtable,array}
\usepackage[round,authoryear]{natbib}
\usepackage[hidelinks]{hyperref}
\setlength{\emergencystretch}{3em}
\newcommand{\resultsfigure}[3]{%
  \begin{figure}[htbp]\centering
  \fbox{\begin{minipage}{0.91\linewidth}\small
  \textit{Figure to be supplied.} #3
  \end{minipage}}
  \caption{#1}\label{#2}\end{figure}}
\begin{document}
\setcounter{chapter}{3}
\chapter{Results}
\label{chap:results}

\section{Introduction}

This chapter reports the observed recognition performance and computational cost of the Shona automatic speech recognition experiments. It addresses two research questions: how curriculum learning affects fine-tuning, and how combining curriculum learning with data augmentation affects in-domain and external-domain recognition. The principal comparisons use Whisper Base on the speaker-disjoint WAXAL v2 partitions described in the methodology. The chapter also reports the completed portion of the matched Whisper Medium extension. Whisper is the pretrained encoder--decoder recogniser adapted in these experiments \citep{Radford2023}.

The presentation distinguishes differences in measured performance from explanations of those differences. It identifies which conditions have lower error, whether word- and character-level outcomes agree, how results vary across seeds and datasets, and what computation was recorded. Explanations involving regularisation, acoustic mismatch, representation learning, or the linguistic properties of Shona are reserved for the discussion chapter. In particular, a lower point estimate alone is not described as proof of a general improvement.

Results reflect completed metric and prediction artifacts available on 13 October 2026. The scope is the two research questions formalised in the methodology, including their model-capacity extension. Prepared follow-up experiments without completed evaluations are not assigned outcomes. Earlier experiments using different partitions are reported only where they clarify the evaluation baseline; they are not pooled with the controlled results.

\section{Reading and interpreting the reported measurements}
\label{sec:results-conventions}

\subsection{Recognition error, loss, and computational cost}

Word error rate (WER) measures the minimum word substitutions ($S$), deletions ($D$), and insertions ($I$) needed to align a generated transcription with its reference, divided by the number of reference words ($N$):
\begin{equation}
  \mathrm{WER}=100\frac{S+D+I}{N}.
\end{equation}
Character error rate (CER) applies the same edit-count calculation to characters. A substitution replaces a reference unit, a deletion omits one, and an insertion adds a unit absent from the reference. All WER and CER values in this chapter are percentages, and lower values indicate fewer transcription errors. Corpus-level scores pool edit counts and reference lengths across recordings; they are not unweighted averages of utterance-level percentages. Because insertions are included in the numerator, an error rate can exceed 100\%.

WER and CER measure different aspects of the output. An incorrectly transcribed word counts as a word error even if only one character differs. CER can therefore be much lower than WER without either metric being incorrect. Neither metric directly establishes intelligibility, semantic adequacy, or the linguistic cause of an error.

Validation loss is the mean cross-entropy recorded when the model predicts reference tokens using preceding reference tokens as context. This procedure is called \emph{teacher forcing}. A token is a unit in the model's text vocabulary and need not be a whole word. Lower loss indicates greater probability assigned to the reference sequence under that evaluation procedure. WER and CER instead score freely generated text. Consequently, lower loss does not necessarily mean lower recognition error. The loss values reported here are dimensionless optimisation diagnostics, not percentages.

Training runtime is elapsed wall-clock time recorded by the training loop, including its scheduled validation evaluations. It is not integrated GPU compute time or a measure of electrical energy. GPU refers to the graphics processing unit used to accelerate training. Peak allocated memory is the largest observed volume of active tensor storage; peak reserved memory also includes capacity retained by the memory allocator. These quantities are distinguished from the GPU's advertised memory capacity and from whole-device memory consumption.

\subsection{Aggregation, uncertainty, and metric provenance}

For conditions with three varying seed outcomes, tables report the arithmetic mean and sample standard deviation (SD) over seeds 42, 43, and 44:
\begin{equation}
 \bar{x}=\frac{1}{3}\sum_{j=1}^{3}x_j,
 \qquad s=\sqrt{\frac{\sum_{j=1}^{3}(x_j-\bar{x})^2}{2}}.
\end{equation}
SD describes the observed spread over these runs. It is not a confidence interval and does not incorporate uncertainty from choosing different speakers or reconstructing the dataset split. A seed changes the pseudorandom sequence used by stochastic operations; it does not provide a new evaluation population. Three seeds permit inspection of variability but give only a limited view of its distribution.

The strict SNR and joint SNR--duration conditions, C2 and C3, have identical saved validation and test prediction hashes across their three seed-labelled exports. Their accuracy is consequently displayed as a single value rather than as a mean with an apparently reassuring zero SD. These records demonstrate prediction repeatability under the deterministic ordering; they do not supply three distinct recognition outcomes. The earlier waveform aggregation reused one C3 prediction export in its clean-curriculum cell. Later C3 seed-44 training and evaluation records are present in the current evidence set and reproduce the same predictions. Only one corrected-FLEURS export exists for C3. Its external-domain score is also reported as a single value. Single-seed C4, C4R, C7, and Medium SortaGrad results are similarly shown without SD, but represent a different limitation: incomplete replication.

Unless explicitly labelled otherwise, recognition tables use the saved item-level prediction summaries, which contain the generated hypotheses' word edit counts and CER. Loss and runtime come from training-result artifacts. Training-loop generation and the separately executed prediction export sometimes give slightly different WER values. For example, Base C0 seed 42 has a final training-loop test WER of 31.3400\%, whereas the item-level export gives 31.3560\%. These are distinct evaluation passes. They are not averaged together, and one measurement source is used consistently within each recognition comparison. Learning-curve values retain their training-loop source.

An improvement is reported in \emph{percentage points} (pp), defined as control error minus treatment error. Thus a positive improvement denotes lower treatment error, and a change from 31.04\% to 30.20\% is approximately 0.84 pp. This is different from a relative percentage reduction. Calculations use unrounded source values; displayed table entries are rounded afterwards.

Where available, paired bootstrap intervals compare two fixed checkpoints on the same evaluation population. The procedure repeatedly resamples groups of recordings and recalculates their difference in corpus WER. WAXAL intervals group recordings by speaker; corrected FLEURS intervals group them by sentence prompt because speaker identifiers are unavailable. The archived analyses use 10,000 replicates. Such intervals describe evaluation-sample uncertainty conditional on the compared checkpoints, rather than training-seed uncertainty. The initial multi-condition curriculum screen has unadjusted intervals: no claim of family-wise statistical significance is made. No new significance claim is inferred from SD overlap or from a three-seed average.

\section{Baseline and evaluation-protocol results}

\subsection{Earlier partitions and the controlled baseline}

Speaker-disjoint partitions contain no speaker in more than one split. The final WAXAL v2 protocol contains 13,807 training, 1,715 validation, and 1,671 test recordings, with 113, 24, and 24 speakers, respectively. Validation is used for checkpoint selection; test evaluation describes the selected checkpoint. The baseline uses shuffled training order without augmentation. A checkpoint is a saved model state; selection uses the lowest recorded validation WER, which need not occur at the last training epoch.

Table~\ref{tab:results-protocol} separates the earlier partition diagnostics from the controlled v2 baseline. The first speaker-disjoint partition, v1, had substantial concentration: one speaker contributed approximately 51.1\% of validation recordings and another approximately 35.1\% of test recordings. Version 2 increased held-out speaker coverage and restricted each held-out speaker's share to at most 20\% of both recordings and duration. These protocols therefore assess different populations and are not interchangeable baseline runs.

\begin{table}[htbp]\centering\small
\caption{Partition diagnostics for Base, seed 42. All four error columns use item-level prediction summaries.}\label{tab:results-protocol}
\begin{tabular}{lrrrr}\toprule
Partition & Val. WER & Test WER & Val. CER & Test CER\\\midrule
Original, speaker overlap & 29.30 & 28.72 & 6.69 & 6.23\\
Speaker-disjoint v1 & 30.19 & 28.26 & 6.06 & 6.00\\
Speaker-disjoint v2 & 32.52 & 31.36 & 7.08 & 7.01\\\bottomrule
\end{tabular}\end{table}

The v2 seed-42 baseline has higher validation and test error than either earlier diagnostic. However, disjointness, speaker composition, partition membership, and some evaluation scheduling differ across these records. The observed differences therefore do not isolate a numerical effect attributable solely to speaker overlap. The earlier v1 test result of 28.26\% is not used as the control for subsequent v2 curricula.

Across the three v2 random-order runs, validation WER is $32.91\pm0.48$\% and test WER is $31.04\pm0.29$\%. Validation and test CER are $7.32\pm0.24$\% and $6.82\pm0.17$\%, respectively. These values are the clean Base reference throughout the chapter.

\subsection{Pretrained Base without study-specific fine-tuning}

The retrospective zero-shot evaluation measures the pretrained Base checkpoint before any Shona updates by this study. It does not establish that Shona was absent from the model's upstream pretraining. Table~\ref{tab:results-zero} compares this reference with the clean fine-tuned baseline under the same saved evaluation manifests. FLEURS, the Few-shot Learning Evaluation of Universal Representations of Speech benchmark, supplies the external Shona speech corpus \citep{Conneau2022}. Its corrected local version preserves the evaluation recordings while correcting the text and identifier preparation described in the methodology.

\begin{table}[htbp]\centering\small
\caption{Zero-shot Base and clean fine-tuned Base. Fine-tuned entries are three-seed means; the zero-shot model is evaluated once.}\label{tab:results-zero}
\begin{tabular}{llrrrr}\toprule
Corpus & Model & Val. WER & Test WER & Val. CER & Test CER\\\midrule
WAXAL v2 & Zero-shot & 134.71 & 132.17 & 56.92 & 57.21\\
WAXAL v2 & Fine-tuned C0 & 32.91 & 31.04 & 7.32 & 6.82\\
FLEURS & Zero-shot & 163.43 & 155.37 & 76.84 & 71.81\\
FLEURS & Fine-tuned C0 & 56.94 & 58.31 & 20.19 & 18.93\\\bottomrule
\end{tabular}\end{table}

Fine-tuning is associated with a WAXAL test WER reduction of 101.12 pp and a CER reduction of 50.39 pp relative to the evaluated pretrained checkpoint. The corresponding FLEURS reductions are 97.06 and 52.88 pp. These large differences establish the scale of adaptation under this protocol; the curriculum and augmentation comparisons concern much smaller differences among fine-tuned models. Values above 100\% in the zero-shot rows are valid edit rates, reflecting the inclusion of inserted words as well as substitutions and deletions.

\section{RQ1: Curriculum learning}
\label{sec:results-rq1}

\subsection{Conditions and aggregate recognition outcomes}

Curriculum learning changes the sequence or eligibility of training examples using a difficulty proxy. A static proxy is fixed before training; a dynamic proxy is updated using model behaviour. The conditions are identified consistently as follows. C0 shuffles the data. C1 applies a duration-based SortaGrad schedule, presenting shorter recordings first in the first epoch and shuffling thereafter, following the duration-ordering idea of \citet{Amodei2016}. An epoch is one scheduled pass through the training exposure budget. C2 presents recordings in descending estimated signal-to-noise ratio (SNR), a measure of speech energy relative to estimated background energy. C3 combines duration and inverse-SNR percentile ranks, with easier-ranked items first in each epoch. These acoustic quantities are proxies, not measured linguistic difficulty.

C4 progressively expands eligibility through acoustic tiers; C4R applies a matched random-pacing control. Pacing changes which recordings are eligible at a given stage, and may repeat eligible recordings to maintain the exposure budget. C5 uses sequence-to-sequence (S2S) negative log-likelihood normalised by audio duration to order later epochs. Negative log-likelihood measures how improbable the reference tokens are under the current model. C6 uses generated utterance-level WER as its updated difficulty score. C7 combines the fixed acoustic ranking with a model-derived S2S ranking. The historical configuration names containing ``margin'' are not interpreted as a separate measured classification margin.

\begin{table}[htbp]\centering\scriptsize
\caption{Base curriculum results on WAXAL v2. Varying three-seed outcomes are mean $\pm$ sample SD. C2/C3 have prediction-identical seed-labelled exports; C4/C4R/C7 have one seed.}\label{tab:results-rq1}
\setlength{\tabcolsep}{4pt}
\begin{tabular}{lrrrr}\toprule
Condition & Val. WER & Test WER & Val. CER & Test CER\\\midrule
% BEGIN GENERATED rq1
C0: Random & $32.91 \pm 0.48$ & $31.04 \pm 0.29$ & $7.32 \pm 0.24$ & $6.82 \pm 0.17$ \\
C1: SortaGrad & $32.60 \pm 0.11$ & $30.96 \pm 0.06$ & $7.26 \pm 0.06$ & $6.89 \pm 0.12$ \\
C2: Strict SNR & $33.12$ & $31.53$ & $7.36$ & $7.46$ \\
C3: SNR--duration & $32.84$ & $30.46$ & $7.13$ & $6.47$ \\
C4: Acoustic pacing & $34.47$ & $32.59$ & $7.81$ & $7.47$ \\
C4R: Random pacing & $34.39$ & $32.46$ & $7.83$ & $7.45$ \\
C5: S2S loss & $33.91 \pm 0.26$ & $31.78 \pm 0.73$ & $7.92 \pm 0.01$ & $7.37 \pm 0.68$ \\
C6: Utterance WER & $33.92 \pm 0.33$ & $31.98 \pm 0.28$ & $7.47 \pm 0.22$ & $6.93 \pm 0.10$ \\
C7: Acoustic--S2S & $34.52$ & $32.19$ & $8.41$ & $7.57$ \\
% END GENERATED rq1
\bottomrule\end{tabular}\end{table}

Among the curricula with varying three-seed outcomes, SortaGrad has the lowest mean validation and test WER: 32.60\% and 30.96\%. Relative to C0, these correspond to reductions of 0.30 and 0.09 pp. Validation CER is also lower by 0.05 pp, but test CER is higher by 0.07 pp. Thus its small WER advantage is not accompanied by improvement in every metric. Its test WER values range from 30.90\% to 31.02\%, compared with 30.79\% to 31.36\% for C0. This describes the observed runs, not a general estimate of robustness to arbitrary seeds.

Strict SNR ordering, C2, gives validation/test WER of 33.12/31.53\% and CER of 7.36/7.46\%. All four errors exceed the corresponding C0 means. Joint SNR--duration ordering, C3, has the lowest clean-curriculum test WER and CER, 30.46\% and 6.47\%, reducing the C0 means by 0.59 and 0.35 pp. Its validation WER, 32.84\%, is only 0.07 pp below the C0 mean and remains above SortaGrad's 32.60\%. The validation and test rankings therefore differ. Furthermore, the identical C3 predictions do not supply three independent confirmations of that ranking.

The S2S and utterance-WER curricula have mean test WER of 31.78\% and 31.98\%, respectively, compared with 31.04\% for random ordering. Their validation WER and both CER means also exceed C0. C5 has the largest test-WER SD among these repeated curriculum conditions, 0.73 pp. Its seed-42 test WER is slightly below the matched C0 value, but the direction does not persist at seeds 43 and 44. A favourable result from that one seed would therefore misrepresent the aggregated outcome.

The single-seed pacing conditions produce test WER of 32.59\% for C4 and 32.46\% for C4R, both above the matched C0 seed-42 value of 31.36\%. Their difference is only 0.13 pp. C7 gives test WER of 32.19\% and the highest validation CER in the curriculum table, 8.41\%. These are observations from the initial screen; absent additional seeds, they are not assigned a run-to-run SD.

\subsection{Paired evidence from the initial screen}

Table~\ref{tab:results-bootstrap} provides the archived seed-42 paired comparisons. Positive values favour the curriculum. The validation intervals for C2 and C4--C7 lie below zero, whereas the SortaGrad interval spans zero. C3 has opposite validation and test directions: its validation estimate favours C0, while its test estimate favours C3. C4 and C4R have wholly negative test intervals. The remaining test intervals, other than C3, span zero. An interval spanning zero leaves the direction unresolved at this interval level; it is not evidence of equivalence.

\begin{table}[htbp]\centering\small
\caption{Seed-42 WER improvement against C0 in pp, with unadjusted 95\% speaker-clustered bootstrap intervals. These fixed-checkpoint intervals are distinct from across-seed SD.}\label{tab:results-bootstrap}
\begin{tabular}{lrr}\toprule
Condition & Validation improvement [CI] & Test improvement [CI]\\\midrule
% BEGIN GENERATED screen
C1: SortaGrad & $-0.07\;[-0.64, +0.27]$ & $+0.34\;[-0.55, +1.00]$ \\
C2: Strict SNR & $-0.60\;[-1.07, -0.19]$ & $-0.18\;[-1.58, +1.01]$ \\
C3: SNR--duration & $-0.32\;[-0.75, -0.02]$ & $+0.90\;[+0.31, +1.39]$ \\
C4: Acoustic pacing & $-1.95\;[-2.76, -1.27]$ & $-1.23\;[-1.77, -0.74]$ \\
C4R: Random pacing & $-1.87\;[-2.78, -1.17]$ & $-1.10\;[-1.61, -0.67]$ \\
C5: S2S loss & $-1.32\;[-2.15, -0.59]$ & $+0.04\;[-0.94, +0.92]$ \\
C6: Utterance WER & $-1.51\;[-2.18, -0.95]$ & $-0.42\;[-1.41, +0.52]$ \\
C7: Acoustic--S2S & $-2.00\;[-3.29, -1.08]$ & $-0.84\;[-2.81, +0.71]$ \\
% END GENERATED screen
\bottomrule\end{tabular}\end{table}

The curriculum with the strongest seed-42 WAXAL test estimate, C3, was subsequently carried into the waveform experiment. Since WAXAL test performance contributed to that historical selection, its favourable test result is selection-stage evidence. It is not presented as a previously untouched confirmation of a validation-only choice. The external-domain results below provide a separate cross-corpus observation.

\resultsfigure{Curriculum error differences on WAXAL v2.}{fig:results-curricula}{Create two aligned dot plots, one for validation WER and one for test WER. Plot C0--C7 including C4R, show individual seed points and the mean for varying runs, and distinguish deterministic C2/C3 values and single-seed conditions by symbol. If bootstrap intervals are included, place them in a separate seed-42 panel labelled ``control minus treatment, pp''; do not mix them with seed SD. Use Table~\ref{tab:results-rq1} and Table~\ref{tab:results-bootstrap}.}

\section{RQ2: Data augmentation and curriculum combinations}
\label{sec:results-rq2}

Data augmentation modifies the training inputs while retaining the reference transcription. Evaluation in this section uses unaugmented WAXAL v2 speech. Two separate experiment families were completed: composite waveform augmentation with random or C3 ordering, and spectrogram masking with random or SortaGrad ordering. The curriculum differs between families, so an augmentation effect is measured against the clean control with the same ordering.

\subsection{Composite waveform augmentation}

The waveform policy applies tempo perturbation, simulated reverberation, and additive noise according to the fixed probabilities and severity settings in the methodology. A materialised view is an audio variant generated in advance and reused during training. A0 denotes the clean random control C0; A1 denotes the clean C3 control; A2 and A3 apply the waveform policy to random and C3 ordering, respectively. Table~\ref{tab:results-wave} reports the four cells of this factorial design, meaning that both ordering conditions are observed with and without augmentation.

\begin{table}[htbp]\centering\scriptsize
\caption{Waveform factorial on WAXAL v2. C0 and C3 are the reused clean A0 and A1 controls; their inclusion does not add new training runs.}\label{tab:results-wave}
\setlength{\tabcolsep}{4pt}
\begin{tabular}{lrrrr}\toprule
Condition & Val. WER & Test WER & Val. CER & Test CER\\\midrule
% BEGIN GENERATED waveform
C0: Random & $32.91 \pm 0.48$ & $31.04 \pm 0.29$ & $7.32 \pm 0.24$ & $6.82 \pm 0.17$ \\
C3: SNR--duration & $32.84$ & $30.46$ & $7.13$ & $6.47$ \\
A2: Waveform, random & $33.06 \pm 0.24$ & $31.30 \pm 0.28$ & $8.01 \pm 0.47$ & $7.12 \pm 0.22$ \\
A3: Waveform, C3 & $33.04 \pm 0.34$ & $31.37 \pm 0.27$ & $7.72 \pm 0.13$ & $7.10 \pm 0.15$ \\
% END GENERATED waveform
\bottomrule\end{tabular}\end{table}

With random ordering, waveform augmentation increases mean validation WER from 32.91\% to 33.06\% and test WER from 31.04\% to 31.30\%. Validation and test CER increase by 0.69 and 0.30 pp, respectively. Accordingly, all four aggregate errors are higher than the clean random control.

With C3 ordering, waveform augmentation increases validation WER from 32.84\% to 33.04\% and test WER from 30.46\% to 31.37\%. The corresponding CER changes are increases of 0.59 and 0.64 pp. Within the augmented conditions, C3 has a validation-WER advantage of only 0.03 pp over random ordering, while its test WER is 0.07 pp higher. The clean C3 test-ranking advantage is therefore absent in the augmented pair.

The descriptive interaction asks whether the ordering advantage changes when augmentation is applied. Defining it as $(\mathrm{A2}-\mathrm{A3})-(\mathrm{A0}-\mathrm{A1})$ on WER, with positive values denoting an increased curriculum advantage, gives $-0.04$ pp on validation and $-0.66$ pp on test. These are differences of observed cell values, not formal interaction-test results. In particular, the A1 cell shares deterministic predictions across its aggregate positions.

The recorded test error decomposition provides a more specific description of A2. Relative to the number of reference words, mean insertions increase from 4.543\% to 4.829\% and substitutions from 22.694\% to 22.771\%, while deletions decrease from 3.807\% to 3.697\%. Thus a reduction in one error component coexists with a higher total WER; no causal explanation for that redistribution is inferred here.

\subsection{SpecAugment policy definitions}

SpecAugment masks parts of the input time--frequency representation \citep{Park2019}. A frequency mask obscures consecutive frequency bands, and a time mask obscures consecutive frames. The completed experiments use the LibriSpeech Basic (LB) and LibriSpeech Double (LD) masking parameters without time warping. The local adaptations may also be termed WAXAL Basic (WB) and WAXAL Double (WD), respectively. The labels LB and LD are used here to maintain a consistent correspondence with the recorded experiment configurations.

\begin{table}[htbp]\centering\small
\caption{Implemented SpecAugment masks. Widths are maxima, not fixed numbers of masked positions; overlaps and zero-width draws are possible.}\label{tab:results-masks}
\begin{tabular}{lrrrrr}\toprule
Policy & Freq. width & Freq. masks & Time width & Time masks & Time ratio\\\midrule
LB, no warp & 27 & 1 & 100 & 1 & 1.0\\
LD, no warp & 27 & 2 & 100 & 2 & 1.0\\\bottomrule
\end{tabular}\end{table}

Frequency width is measured in input bands and time width in feature frames. The time ratio limits the width of each time-mask attempt relative to the valid recording frames; it is not the number of masks. LD doubles the mask counts while preserving the individual width limits. It does not set the time ratio to 2.0. Masks are applied during training, with waveform augmentation disabled, and are absent during evaluation.

\subsection{SpecAugment recognition outcomes}

\begin{table}[htbp]\centering\scriptsize
\caption{SpecAugment on WAXAL v2, mean $\pm$ sample SD over three seeds. C0 and C1 supply the clean random and clean SortaGrad controls.}\label{tab:results-specaug}
\setlength{\tabcolsep}{4pt}
\begin{tabular}{lrrrr}\toprule
Condition & Val. WER & Test WER & Val. CER & Test CER\\\midrule
% BEGIN GENERATED specaugment
C0: Random & $32.91 \pm 0.48$ & $31.04 \pm 0.29$ & $7.32 \pm 0.24$ & $6.82 \pm 0.17$ \\
C1: SortaGrad & $32.60 \pm 0.11$ & $30.96 \pm 0.06$ & $7.26 \pm 0.06$ & $6.89 \pm 0.12$ \\
S2: LB, random & $31.76 \pm 0.16$ & $30.20 \pm 0.35$ & $6.87 \pm 0.11$ & $6.71 \pm 0.45$ \\
S3: LB, SortaGrad & $32.20 \pm 0.26$ & $30.43 \pm 0.22$ & $7.22 \pm 0.26$ & $6.89 \pm 0.26$ \\
S4: LD, random & $32.46 \pm 0.28$ & $30.38 \pm 0.19$ & $7.54 \pm 0.33$ & $6.93 \pm 0.28$ \\
S5: LD, SortaGrad & $32.41 \pm 0.08$ & $30.69 \pm 0.41$ & $7.30 \pm 0.07$ & $7.09 \pm 0.30$ \\
% END GENERATED specaugment
\bottomrule\end{tabular}\end{table}

LB masking with random ordering, S2, gives the lowest mean WAXAL validation and test WER among the completed Base conditions: $31.76\pm0.16$\% and $30.20\pm0.35$\%. Relative to C0, the improvements are 1.14 and 0.84 pp. Mean validation CER decreases by 0.45 pp and test CER by 0.11 pp. All three S2 seeds have lower validation and test WER than their seed-matched C0 controls. Nevertheless, seed-42 test CER rises from 7.01\% to 7.22\%, so the mean CER improvement is not uniform across seeds. C3 retains a lower test CER of 6.47\% than S2's mean of 6.71\%; the condition with the lowest mean test WER is therefore not best on every endpoint.

LB masking with SortaGrad, S3, improves mean validation and test WER over clean SortaGrad by 0.41 and 0.53 pp, respectively. Its mean test CER is almost unchanged, 6.8853\% compared with 6.8895\%. Relative to LB with random ordering, however, LB with SortaGrad has higher validation WER by 0.44 pp and higher test WER by 0.23 pp. Adding SortaGrad does not improve the observed LB cell means.

LD with random ordering, S4, reduces mean test WER by 0.66 pp relative to C0, but it is 0.18 pp worse than LB with random ordering. Its mean validation and test CER, 7.54\% and 6.93\%, both exceed the clean random means. LD with SortaGrad, S5, reduces test WER by 0.26 pp relative to clean SortaGrad, while increasing test CER by 0.20 pp. It also has higher validation and test WER than the corresponding LB--SortaGrad condition. Doubling mask counts therefore does not improve the mean WAXAL WER within either ordering in these runs.

These comparisons establish an observed in-domain ranking. The means and SDs alone do not establish a population-level treatment effect, and no unreported multiple-comparison adjustment is assumed. In particular, improvements reported for SpecAugment in its original study do not determine whether the present models improve on the external Shona corpus; that question is addressed by their measured FLEURS outcomes.

\resultsfigure{Curriculum and augmentation combinations on WAXAL v2.}{fig:results-factorial}{Create separate panels for waveform and spectrogram augmentation. In the waveform panel, plot clean versus waveform WER for random and C3 ordering. In the spectrogram panel, plot clean, LB, and LD WER for random and SortaGrad ordering. Use validation and test subpanels, individual seed points, and labelled mean lines. Mark the deterministic clean C3 point once. Do not join C3 and SortaGrad as though they were the same curriculum.}

\section{External-domain evaluation on corrected FLEURS}
\label{sec:results-fleurs}

The external evaluation uses 393 validation and 925 test recordings from the corrected Shona FLEURS preparation. These recordings were not used in controlled WAXAL fine-tuning. Out-of-domain here means a different corpus and collection process; it does not assert that every recording was absent from Whisper's upstream pretraining. All conditions in Table~\ref{tab:results-fleurs} are scored against the same corrected references. Conditions without a saved external evaluation are omitted rather than assigned an estimated outcome.

\begin{table}[htbp]\centering\scriptsize
\caption{Corrected FLEURS recognition results. Three-seed conditions use mean $\pm$ sample SD; C3 uses its one available external prediction export.}\label{tab:results-fleurs}
\setlength{\tabcolsep}{4pt}
\begin{tabular}{lrrrr}\toprule
Condition & Val. WER & Test WER & Val. CER & Test CER\\\midrule
% BEGIN GENERATED fleurs
C0: Random & $56.94 \pm 1.25$ & $58.31 \pm 0.59$ & $20.19 \pm 1.66$ & $18.93 \pm 0.26$ \\
C1: SortaGrad & $55.74 \pm 0.31$ & $57.76 \pm 0.05$ & $19.13 \pm 0.88$ & $19.17 \pm 0.32$ \\
C3: SNR--duration & $59.62$ & $59.23$ & $21.39$ & $20.08$ \\
A2: Waveform, random & $56.95 \pm 0.72$ & $59.49 \pm 1.47$ & $19.34 \pm 1.33$ & $20.03 \pm 0.42$ \\
A3: Waveform, C3 & $57.04 \pm 0.94$ & $60.10 \pm 2.04$ & $19.53 \pm 0.96$ & $20.56 \pm 1.60$ \\
S2: LB, random & $57.12 \pm 1.13$ & $59.67 \pm 0.51$ & $19.89 \pm 0.42$ & $20.47 \pm 0.34$ \\
S3: LB, SortaGrad & $57.17 \pm 0.94$ & $59.17 \pm 1.03$ & $19.80 \pm 0.82$ & $20.16 \pm 0.88$ \\
S4: LD, random & $62.10 \pm 0.99$ & $62.83 \pm 1.10$ & $24.48 \pm 0.82$ & $22.21 \pm 0.93$ \\
S5: LD, SortaGrad & $59.80 \pm 1.42$ & $63.03 \pm 1.24$ & $21.54 \pm 0.41$ & $22.55 \pm 1.29$ \\
% END GENERATED fleurs
\bottomrule\end{tabular}\end{table}

All evaluated fine-tuned Base conditions have substantially higher error on FLEURS than on WAXAL. The C0 mean test WER increases from 31.04\% on WAXAL to 58.31\% on FLEURS, a 27.27 pp cross-corpus difference; test CER increases from 6.82\% to 18.93\%. This difference describes the two evaluation populations and reference protocols, rather than isolating a particular source of domain difficulty.

Clean SortaGrad, C1, has the lowest mean FLEURS validation and test WER among the evaluated conditions: 55.74\% and 57.76\%. Its mean test WER is 0.55 pp below C0, but its test CER is 0.24 pp higher. Thus the cross-corpus WER ranking does not yield a corresponding best CER ranking. C3 records test WER of 59.23\%, 0.92 pp above the C0 mean despite its lower WAXAL test WER. The archived prompt-clustered C3-versus-C0 intervals span zero for each matched baseline seed. The observed external ranking therefore supplies no confirmed C3 improvement.

Waveform augmentation gives mean FLEURS test WER of 59.49\% with random ordering and 60.10\% with C3. These exceed the clean random mean by 1.18 and 1.79 pp, respectively; A3 is also 0.87 pp above clean C3. The A2 seed effects against C0 vary in direction: the control-minus-treatment WER differences are $-3.47$, $+0.12$, and $-0.18$ pp. In the saved paired analysis, only the seed-42 interval excludes zero, with $[-7.11,-0.22]$ pp favouring the clean control. This is a fixed-checkpoint comparison, not a three-seed significance test.

The in-domain advantage of LB masking does not carry over to a lower mean FLEURS test WER. S2 gives 59.67\%, 1.36 pp above clean random ordering. S3 gives 59.17\%, 1.41 pp above clean SortaGrad. Both also have higher mean test CER than their corresponding clean controls. LD increases external error further: S4 and S5 reach mean test WER of 62.83\% and 63.03\%, respectively, with test CER of 22.21\% and 22.55\%. The S2 condition that leads the mean WAXAL WER ranking therefore does not lead the FLEURS ranking, and stronger masking does not improve the observed external-domain averages.

The prompt-clustered analyses use 348 distinct test sentence prompts, grouping repeated readings of the same prompt together. This recognises dependence due to shared text but cannot account for unknown speaker dependence. No complete paired significance analysis of the SpecAugment external comparisons is supplied in the evidence used here; their differences are reported descriptively. More generally, the observed cross-corpus differences do not identify whether acoustics, textual content, transcription conventions, or another property accounts for the ranking change.

\resultsfigure{WAXAL and corrected FLEURS test-error rankings.}{fig:results-domain}{Create paired panels for test WER and test CER, with one point per condition for each corpus. Show C0, C1, C3, A2, A3, and S2--S5. Display individual seed points lightly behind the means; show C3 once. Either use clearly labelled separate corpus axes or a connected-dot layout with sufficient range. The figure should make the S2 in-domain WER advantage and its absence on FLEURS visible without implying that a cross-corpus difference is a paired treatment effect.}

\section{Learning curves and checkpoint selection}
\label{sec:results-curves}

Table~\ref{tab:results-curves} reports representative logged validation WER and loss at the three epoch boundaries. For these v2 runs, the boundaries are steps 2,302, 4,604, and 6,906. A step is a recorded training update position. The available Base C1 seed-43 history is used because its seed-42 training history is not present in the local mirror. These examples document the trajectories without assuming that an absent history can be reconstructed from a final score.

\begin{table}[htbp]\centering\small
\caption{Logged validation WER (\%) / validation loss at epoch boundaries. These are training-loop measurements, not the separately exported recognition scores.}\label{tab:results-curves}
\setlength{\tabcolsep}{4pt}
\begin{tabular}{lrrr}\toprule
Condition & Epoch 1 & Epoch 2 & Epoch 3\\\midrule
% BEGIN GENERATED curves
C0, seed 42 & $37.08/0.4409$ & $33.18/0.3887$ & $32.53/0.3815$ \\
C1, seed 43 & $43.55/0.4471$ & $32.95/0.3916$ & $32.46/0.3835$ \\
C5, seed 42 & $37.08/0.4409$ & $33.82/0.3896$ & $34.48/0.3864$ \\
C6, seed 42 & $37.08/0.4409$ & $34.04/0.3907$ & $35.11/0.3867$ \\
A2, seed 42 & $36.78/0.4386$ & $34.50/0.3869$ & $33.17/0.3774$ \\
S2, seed 42 & $37.36/0.4345$ & $33.33/0.3787$ & $31.69/0.3657$ \\
S4, seed 42 & $37.35/0.4419$ & $34.47/0.3826$ & $32.61/0.3680$ \\
M0, seed 42 & $28.17/0.3016$ & $26.02/0.2726$ & $26.50/0.2774$ \\
M1, seed 42 & $63.08/0.3129$ & $26.42/0.2739$ & $26.27/0.2784$ \\
% END GENERATED curves
\bottomrule\end{tabular}\end{table}

For Base random ordering, validation WER and loss decrease across all three logged epochs for all seeds. In seed 42, WER falls from 37.08\% to 33.18\%, then to 32.53\%. The second improvement is smaller, but the direction remains favourable. It would therefore be inaccurate to describe the epoch-two boundary as a universal onset of deteriorating WER. The available clean SortaGrad histories also show decreasing WER and loss after epoch two, despite higher first-epoch WER than the random baseline.

The dynamic C5 and C6 conditions behave differently. In C5 seed 42, WER rises from 33.82\% at epoch two to 34.48\% at epoch three, while loss falls from 0.3896 to 0.3864. C6 seed 42 similarly rises from 34.04\% to 35.11\% WER while loss falls from 0.3907 to 0.3867. The same divergence in direction occurs in their seed-43 and seed-44 histories. Their selected checkpoints are therefore from epoch two, even though optimisation continues through epoch three. The final recognition tables evaluate those selected checkpoints, not the last epoch merely because it was trained most recently.

All completed Base waveform and SpecAugment runs retain the epoch-three checkpoint, with validation WER decreasing across the three logged boundaries. LB random seed 42 falls from 37.36\% to 33.33\% to 31.69\%, while loss falls from 0.4345 to 0.3787 to 0.3657. Across seeds, the selected LB-random validation loss averages 0.3660 compared with 0.3821 for C0. In this comparison lower loss accompanies lower WAXAL error. By contrast, waveform-random mean loss is also lower than C0, at 0.3782, but its WER and CER means are higher. The relationship between loss and recognition error therefore differs across conditions.

These histories establish the direction of change at recorded evaluation boundaries. Three evaluation points are insufficient to locate an exact within-epoch change point or to establish a smooth curve's curvature. Training-loss logs sampled more frequently also measure a different quantity from validation loss. Any W\&B visualisation should preserve these distinctions, indicate smoothing explicitly, and mark the checkpoint selected by validation WER.

\resultsfigure{Validation WER and loss across training.}{fig:results-curves}{Export unsmoothed epoch-evaluation histories from W\&B or recreate them from the saved trainer states. Use aligned WER and loss panels for clean C0/C1, dynamic C5/C6, and augmentation A2/S2/S4. Show separate seed traces, epoch boundaries, and selected checkpoints. Label WER as a percentage and loss as cross-entropy. If first-epoch values require a large vertical range, provide a labelled inset for epochs two and three instead of silently clipping values. Do not invent a seed-42 C1 trace when its history is unavailable.}

\section{Computational cost and memory observations}
\label{sec:results-resources}

\subsection{Recorded training runtime}

The controlled run manifests identify an NVIDIA GeForce RTX 4090. Table~\ref{tab:results-resources} summarises recorded runtime and selected-checkpoint validation loss. Training-loop runtime already includes scheduled epoch evaluations; adding them again would double-count time. It excludes separately executed final evaluations and prediction exports, and does not provide a complete end-to-end accounting of preprocessing, waveform materialisation, queue waiting, or later external-domain evaluation. Small timing differences consequently should not be interpreted as established hardware-independent speed differences.

\begin{table}[htbp]\centering\small
\caption{Training-loop wall-clock time in minutes and selected-checkpoint validation loss. Runtime summaries include the recorded executions, even when their predictions are identical. $n$ is the number of run records, not necessarily independent recognition outcomes.}\label{tab:results-resources}
\begin{tabular}{lrrr}\toprule
Condition & $n$ & Runtime (min) & Validation loss\\\midrule
% BEGIN GENERATED resources
C0: Random & 3 & $37.17 \pm 0.55$ & $0.3821 \pm 0.0006$ \\
C1: SortaGrad & 3 & $38.27 \pm 1.10$ & $0.3831 \pm 0.0003$ \\
C2: Strict SNR & 3 & $37.76 \pm 0.99$ & $0.3853$ \\
C3: SNR--duration & 3 & $37.66 \pm 0.92$ & $0.3856$ \\
C4: Acoustic pacing & 1 & $37.09$ & $0.4135$ \\
C4R: Random pacing & 1 & $40.21$ & $0.4079$ \\
C5: S2S loss & 3 & $47.84 \pm 16.84$ & $0.3902 \pm 0.0011$ \\
C6: Utterance WER & 3 & $189.90 \pm 38.52$ & $0.3906 \pm 0.0002$ \\
C7: Acoustic--S2S & 1 & $41.69$ & $0.3882$ \\
A2: Waveform, random & 3 & $37.52 \pm 0.25$ & $0.3782 \pm 0.0008$ \\
A3: Waveform, C3 & 3 & $37.09 \pm 0.69$ & $0.3814 \pm 0.0003$ \\
S2: LB, random & 3 & $38.10 \pm 1.06$ & $0.3660 \pm 0.0003$ \\
S3: LB, SortaGrad & 3 & $38.69 \pm 0.75$ & $0.3666 \pm 0.0002$ \\
S4: LD, random & 3 & $37.49 \pm 0.40$ & $0.3683 \pm 0.0004$ \\
S5: LD, SortaGrad & 3 & $38.04 \pm 0.80$ & $0.3693 \pm 0.0006$ \\
M0: Medium, random & 3 & $151.70 \pm 0.54$ & $0.2759 \pm 0.0030$ \\
M1: Medium, SortaGrad & 1 & $163.42$ & $0.2784$ \\
% END GENERATED resources
\bottomrule\end{tabular}\end{table}

Base random ordering averages 37.17 minutes. Most static, waveform, and masking conditions average approximately 37--39 minutes. LB random averages 38.10 minutes, 0.93 minutes above C0, while giving the lower mean WAXAL WER reported earlier. These elapsed-time observations do not by themselves quantify GPU utilisation, energy demand, or the incremental cost of preparing augmentation assets.

C6, which uses generated utterance WER for difficulty scoring, averages 189.90 minutes, approximately 5.11 times the C0 mean. Its individual runtimes are 234.38, 167.19, and 168.13 minutes for seeds 42--44. The seed-42-only comparison is approximately 6.39 times its matched control, explaining why an earlier report gave a larger ratio. Both comparisons show substantially longer recorded training, while the aggregate recognition errors remain above C0.

C5 illustrates why elapsed-time variability requires care. The original seed-42 execution records 67.29 minutes, compared with 38.00 and 38.24 minutes at seeds 43 and 44. Its exact seed-42 repeat records 38.06 minutes while reproducing the original metric trajectory. The original three seed records are retained in the main runtime average of 47.84 minutes; the repeat is not substituted silently or counted as a fourth accuracy observation. This contrast documents variability in execution time, but the available records do not isolate its cause.

\subsection{Peak-memory evidence and its limits}

No comparable full-training peak-memory series for the controlled Base and Medium conditions was found in the saved result summaries. Those quantities are therefore unreported rather than inferred from a 24 GB device specification or a configured free-memory requirement. The available direct memory measurements come from bounded Medium feasibility probes, summarised in Table~\ref{tab:results-memory}.

\begin{table}[htbp]\centering\small
\caption{Medium memory diagnostics, not full-training peaks. GiB denotes $2^{30}$ bytes. The corrected probes use batch size six and each complete two non-skipped optimiser updates.}\label{tab:results-memory}
\begin{tabular}{lrrr}\toprule
Probe & Training allocated & Training reserved & Evaluation reserved\\\midrule
Clean, corrected & 14.235 GiB & 14.941 GiB & 14.994 GiB\\
LD masking & 14.235 GiB & 14.934 GiB & 14.986 GiB\\\bottomrule
\end{tabular}\end{table}

Allocated training peaks are nearly identical in these two probes. The clean probe's highest reported reserved capacity across the displayed phases is 14.994 GiB. An earlier clean probe is excluded because no optimiser update completed and its optimiser state was uninitialised. Its smaller memory reading is consequently not comparable with the corrected probes.

The diagnostics establish that the tested batches and operations completed under the specified memory conditions. They do not measure the largest allocation over a complete training run, memory fragmentation across many hours, or other processes' device use. A full-run memory ranking among curricula or augmentation policies cannot be supplied from these probes. Similarly, neither energy efficiency nor total storage cost is inferred from runtime and diagnostic memory alone.

\resultsfigure{Recognition error and recorded training cost.}{fig:results-cost}{Plot WAXAL test WER against training-loop runtime in minutes, identifying individual seeds and model sizes. Include C0--C7, waveform, masking, and completed Medium conditions; use a logarithmic time axis if needed. Mark the C5 seed-42 rerun separately as a repeatability diagnostic rather than an additional accuracy replicate. A memory panel should be added only after comparable full-run measurements are exported; otherwise retain Table~\ref{tab:results-memory} as a clearly labelled diagnostic table.}

\section{Completed Whisper Medium extension}
\label{sec:results-medium}

The matched capacity extension changes the pretrained variant while retaining the specified v2 data, training budget, and batching contract. The local completed evidence contains three Medium random-order runs and one Medium SortaGrad run, seed 42. Thus it supports a repeated model-size comparison for random ordering and a single-seed ordering comparison within Medium, rather than a completed curriculum-by-augmentation matrix. The planned matrix contains 27 runs, but a configuration or a retained training checkpoint is not evidence of completed final evaluation.

\begin{table}[htbp]\centering\scriptsize
\caption{Completed model-capacity results on WAXAL v2. M0 is Medium random ordering over three seeds; M1 is Medium SortaGrad at seed 42 only.}\label{tab:results-medium}
\setlength{\tabcolsep}{4pt}
\begin{tabular}{lrrrr}\toprule
Condition & Val. WER & Test WER & Val. CER & Test CER\\\midrule
% BEGIN GENERATED medium
C0: Random & $32.91 \pm 0.48$ & $31.04 \pm 0.29$ & $7.32 \pm 0.24$ & $6.82 \pm 0.17$ \\
M0: Medium, random & $25.75 \pm 0.25$ & $23.93 \pm 0.32$ & $5.00 \pm 0.15$ & $4.53 \pm 0.09$ \\
M1: Medium, SortaGrad & $26.26$ & $23.66$ & $5.16$ & $4.49$ \\
% END GENERATED medium
\bottomrule\end{tabular}\end{table}

Medium random ordering reduces mean validation WER from 32.91\% to 25.75\% and test WER from 31.04\% to 23.93\%, improvements of 7.15 and 7.12 pp relative to Base. Mean validation and test CER decrease by 2.31 and 2.29 pp. Every completed Medium random seed has lower WAXAL WER and CER than its matched Base seed. The mean training-loop runtime rises from 37.17 to 151.70 minutes, approximately 4.08 times the Base duration.

Medium SortaGrad seed 42 has validation/test WER of 26.26/23.66\% and CER of 5.16/4.49\%. The corresponding Medium random seed-42 values are 26.03/24.24\% and 5.04/4.62\%. SortaGrad therefore improves the seed-matched test WER by 0.58 pp and test CER by 0.13 pp, but worsens validation WER by 0.23 pp and validation CER by 0.12 pp. Comparing this one treatment run only with the three-seed random mean would obscure the matched comparison and the incomplete replication.

Medium random seed 42 selects epoch two, whereas seeds 43 and 44 select epoch three. Medium SortaGrad seed 42 selects epoch three; its first-epoch validation WER is 63.08\%, then falls to 26.42\% and 26.27\%. This trajectory differs markedly from the first-epoch random-order result, while its selected-checkpoint test WER is lower. An early-epoch ranking consequently does not determine the final checkpoint ranking in this example.

No completed Medium FLEURS evaluation is available in the inspected artifacts, so cross-domain benefits of the larger model cannot be reported here. Likewise, the available Medium results do not establish whether its response to waveform augmentation or masking matches Base. The comparison concerns the pretrained variants as a whole: architecture and pretrained weights change together, so the observed improvement cannot be attributed uniquely to parameter count.

\section{Findings in relation to the research questions}

For RQ1, the observed effect depends on the curriculum, metric, and evaluation population. Clean SortaGrad gives small mean WAXAL WER reductions, with a slight increase in test CER. Deterministic SNR--duration ordering gives the lowest clean-curriculum WAXAL test WER and CER, but its validation ranking is less favourable and its external test WER exceeds the random baseline. Strict SNR ordering, the dynamic curricula, and the single-seed pacing and hybrid conditions do not improve the complete set of reported WAXAL endpoints over their controls. No evaluated clean curriculum leads across all reported validation, test, word, character, and external-domain measures.

For RQ2, the composite waveform policy increases mean WAXAL WER and CER under both studied orderings. LB masking with random ordering yields the lowest mean WAXAL WER among the completed Base conditions, with a test reduction of 0.84 pp relative to clean random ordering. Adding SortaGrad or doubling mask counts does not improve that mean WAXAL WER result. The masking conditions nevertheless have higher mean FLEURS test WER and CER than their corresponding clean controls. The observed in-domain masking advantage is therefore not an observed external-domain advantage under the evaluated protocol.

Loss, runtime, and memory add distinct information. Lower validation loss sometimes accompanies worse generated-text error; dynamic utterance-WER ordering records substantially longer training without a lower mean error; and full-run memory comparisons remain unavailable. The completed Medium random-order runs show lower WAXAL error at approximately four times the Base training-loop duration, while their external performance and most treatment combinations remain unevaluated. These observations define the evidence to be explained in the discussion chapter, without assigning unmeasured mechanisms to the differences.

\clearpage
\section{Per-seed WAXAL recognition record}
\label{sec:results-perseed}

Table~\ref{tab:results-perseed} retains the individual completed-run values underlying the main aggregates. It excludes the C5 exact rerun from the replication count and does not count shared clean controls again under their factorial aliases. C2 and C3 rows reproduce identical prediction exports as described in Section~\ref{sec:results-conventions}. M1 contains only seed 42. The table provides the seed detail within the chapter without requiring access to the experiment dashboard.

\small
\begin{longtable}{lrrrrr}
\caption{Individual WAXAL v2 results (\%). Condition codes are defined in the preceding sections.}\label{tab:results-perseed}\\
\toprule Condition & Seed & Val. WER & Test WER & Val. CER & Test CER\\\midrule
\endfirsthead
\toprule Condition & Seed & Val. WER & Test WER & Val. CER & Test CER\\\midrule
\endhead
\midrule\multicolumn{6}{r}{Continued on next page}\\\endfoot
\bottomrule\endlastfoot
% BEGIN GENERATED perseed
C0 & 42 & 32.52 & 31.36 & 7.08 & 7.01 \\
C0 & 43 & 33.44 & 30.79 & 7.56 & 6.68 \\
C0 & 44 & 32.76 & 30.99 & 7.30 & 6.77 \\
C1 & 42 & 32.59 & 31.02 & 7.21 & 6.79 \\
C1 & 43 & 32.50 & 30.95 & 7.26 & 7.02 \\
C1 & 44 & 32.73 & 30.90 & 7.32 & 6.85 \\
C2 & 42 & 33.12 & 31.53 & 7.36 & 7.46 \\
C2 & 43 & 33.12 & 31.53 & 7.36 & 7.46 \\
C2 & 44 & 33.12 & 31.53 & 7.36 & 7.46 \\
C3 & 42 & 32.84 & 30.46 & 7.13 & 6.47 \\
C3 & 43 & 32.84 & 30.46 & 7.13 & 6.47 \\
C3 & 44 & 32.84 & 30.46 & 7.13 & 6.47 \\
C4 & 42 & 34.47 & 32.59 & 7.81 & 7.47 \\
C4R & 42 & 34.39 & 32.46 & 7.83 & 7.45 \\
C5 & 42 & 33.84 & 31.32 & 7.92 & 6.96 \\
C5 & 43 & 33.69 & 32.62 & 7.91 & 8.15 \\
C5 & 44 & 34.20 & 31.41 & 7.93 & 6.99 \\
C6 & 42 & 34.03 & 31.77 & 7.68 & 6.81 \\
C6 & 43 & 33.55 & 32.30 & 7.23 & 6.99 \\
C6 & 44 & 34.19 & 31.87 & 7.48 & 6.97 \\
C7 & 42 & 34.52 & 32.19 & 8.41 & 7.57 \\
A2 & 42 & 33.12 & 31.62 & 8.54 & 7.36 \\
A2 & 43 & 33.27 & 31.12 & 7.81 & 7.06 \\
A2 & 44 & 32.79 & 31.15 & 7.68 & 6.94 \\
A3 & 42 & 33.07 & 31.36 & 7.73 & 7.21 \\
A3 & 43 & 33.35 & 31.10 & 7.84 & 6.93 \\
A3 & 44 & 32.68 & 31.65 & 7.58 & 7.16 \\
S2 & 42 & 31.71 & 30.57 & 6.79 & 7.22 \\
S2 & 43 & 31.94 & 30.16 & 7.00 & 6.53 \\
S2 & 44 & 31.64 & 29.87 & 6.82 & 6.38 \\
S3 & 42 & 32.00 & 30.61 & 7.04 & 7.18 \\
S3 & 43 & 32.10 & 30.18 & 7.10 & 6.67 \\
S3 & 44 & 32.49 & 30.49 & 7.51 & 6.81 \\
S4 & 42 & 32.76 & 30.51 & 7.92 & 7.16 \\
S4 & 43 & 32.20 & 30.17 & 7.38 & 6.62 \\
S4 & 44 & 32.41 & 30.48 & 7.32 & 7.00 \\
S5 & 42 & 32.44 & 30.86 & 7.29 & 7.32 \\
S5 & 43 & 32.47 & 30.23 & 7.38 & 6.75 \\
S5 & 44 & 32.33 & 30.99 & 7.23 & 7.18 \\
M0 & 42 & 26.03 & 24.24 & 5.04 & 4.62 \\
M0 & 43 & 25.68 & 23.59 & 5.13 & 4.43 \\
M0 & 44 & 25.55 & 23.95 & 4.84 & 4.55 \\
M1 & 42 & 26.26 & 23.66 & 5.16 & 4.49 \\
% END GENERATED perseed
\end{longtable}
\normalsize

\begin{thebibliography}{9}
\bibitem[Amodei et~al.(2016)]{Amodei2016}
Amodei, D., et~al. (2016). Deep Speech 2: End-to-End Speech Recognition in English and Mandarin. \emph{Proceedings of ICML}, 48, 173--182. \url{https://proceedings.mlr.press/v48/amodei16.html}.
\bibitem[Conneau et~al.(2022)]{Conneau2022}
Conneau, A., et~al. (2022). FLEURS: Few-shot Learning Evaluation of Universal Representations of Speech. arXiv:2205.12446. \url{https://arxiv.org/abs/2205.12446}.
\bibitem[Park et~al.(2019)]{Park2019}
Park, D.~S., Chan, W., Zhang, Y., Chiu, C.-C., Zoph, B., Cubuk, E.~D., and Le, Q.~V. (2019). SpecAugment: A Simple Data Augmentation Method for Automatic Speech Recognition. \emph{Proceedings of Interspeech}, 2613--2617. \url{https://doi.org/10.21437/Interspeech.2019-2680}.
\bibitem[Radford et~al.(2023)]{Radford2023}
Radford, A., Kim, J.~W., Xu, T., Brockman, G., McLeavey, C., and Sutskever, I. (2023). Robust Speech Recognition via Large-Scale Weak Supervision. \emph{Proceedings of ICML}, 202, 28492--28518. \url{https://proceedings.mlr.press/v202/radford23a.html}.
\end{thebibliography}
\end{document}
